\documentclass[11pt,a4paper]{article}

\usepackage[margin=1in]{geometry}
\usepackage[utf8]{inputenc}
\usepackage[T1]{fontenc}
\usepackage{amsmath,amssymb}
\usepackage{hyperref}
\usepackage{authblk}

\hypersetup{
    colorlinks=true,
    linkcolor=blue,
    citecolor=blue,
    urlcolor=blue,
    pdftitle={Open-Source Spatiotemporal Traffic Analysis at Scale},
    pdfauthor={Firman Hadi}
}

\title{\textbf{Open-Source Spatiotemporal Traffic Congestion Analysis at City Scale: A Reproducible Python Pipeline Combining PySAL, OSMnx, and Multilevel Modeling}}

\author[1]{Firman Hadi\thanks{Corresponding author: firmanhadi21@lecturer.undip.ac.id}}
\author[2]{Fauzan Abdullah}
\author[1]{Yasser Wahyuddin}
\author[1]{L.M. Sabri}

\affil[1]{Department of Geodetic Engineering, Universitas Diponegoro, Semarang, Indonesia}
\affil[2]{Department of Statistics, Faculty of Sciences and Mathematics, Universitas Diponegoro, Semarang, Indonesia}

\date{}

\begin{document}

\maketitle

\noindent
Urban traffic congestion costs Southeast Asian economies 2--5\% of GDP annually, yet the analytical methods needed to disentangle its spatial and temporal drivers---exploratory spatial statistics, network topology analysis, multilevel variance decomposition---have traditionally required proprietary GIS platforms whose licensing fees are prohibitive for the developing-country cities that need them most. This extended abstract presents \texttt{traffic-congestion-pipeline}, an open-source Python package (v0.4.0, MIT license, \texttt{pip install traffic-congestion-pipeline}) that delivers a complete, reproducible workflow for large-scale spatiotemporal traffic research using exclusively FOSS4G tools. The package, its source code, and full API documentation are publicly available at \url{https://github.com/firmanhadi21/traffic-analyses} and \url{https://firmanhadi21.github.io/traffic-analyses/}.

\medskip
The pipeline integrates five mature open-source components into a unified eleven-command CLI and Python API. (1)~PySAL's \texttt{esda} and \texttt{libpysal} compute Global Moran's~I and Local Indicators of Spatial Association (LISA) with K-nearest-neighbor spatial weights for hotspot detection. (2)~PySAL's \texttt{giddy} module fits classic Markov and Spatial Markov transition models to quantify how congestion hotspots persist over time and whether their transitions depend on neighboring segments' states (spatial contagion). (3)~OSMnx downloads street network graphs and computes betweenness centrality and graph-based capacity-drop detection, linking network topology to congestion. (4)~\texttt{statsmodels.mixedlm} fits nested mixed-effects models---null, temporal, and full---that rigorously partition within-segment (temporal) and between-segment (spatial) variance using absolute speed (km/h) rather than the normalized jam factor, avoiding circularity from free-flow speed normalization. (5)~Uber H3 hexagonal aggregation re-runs spatial autocorrelation tests at multiple resolutions (6--9) to check robustness against the modifiable areal unit problem (MAUP). Each component maps to a dedicated CLI command (\texttt{traffic-pipeline multilevel}, \texttt{markov}, \texttt{speed-validation}, \texttt{h3-robustness}, etc.), allowing researchers to execute any stage independently or chain the full workflow from automated data collection to publication-ready figures.

We apply this pipeline to over 316 million traffic observations collected at 15-minute intervals from the HERE Traffic API across three Indonesian cities over 12 months (March 2025--February 2026): Jakarta (14,912 road segments, population 10.5~million), Bandung (2,918 segments, 2.5~million), and Semarang (1,047 segments, 1.8~million). These cities span a 20$\times$ population range and feature high motorcycle mode shares, representing traffic dynamics markedly different from car-dominated Western cities.

\medskip
Multilevel variance decomposition shows that 88--89\% of total speed variance lies between segments (ICC), reflecting road design differences. Within segments, time-of-day explains 51--64\% of speed fluctuations, while betweenness centrality adds less than 1\% explanatory power beyond road type. ANOVA across four speed metrics confirms this temporal dominance is not a normalization artifact: $\eta^2$ = 8--10\% for jam factor and speed reduction, 5--6\% for absolute speed, and effectively zero for free-flow speed. Centrality correlations with absolute current speed are near zero ($R^2 < 0.003$); moderate centrality--jam factor correlations ($R^2$ = 0.06--0.14) are mediated entirely through free-flow speed. Evening peak congestion exceeds daily averages by approximately 40\% across all three cities.

Global Moran's~I is significant for all cities ($I = 0.29$--$0.31$, all $p < 0.001$). LISA identifies local clusters in 34--38\% of segments, with 18--25\% surviving FDR correction. LISA Markov analysis reveals an inverse relationship between city size and hotspot persistence: Semarang retains hotspots at higher probability than Jakarta, yet no segment persists across all eight daily periods. Spatial Markov testing detects significant contagion in Bandung ($\chi^2 = 8.43$, $p = 0.004$) and Semarang ($\chi^2 = 6.48$, $p = 0.011$), but not in Jakarta, where denser network topology dissipates spillovers. H3 robustness analysis confirms null spatial autocorrelation at neighbourhood scales for Bandung and Semarang, while revealing a weak signal for Jakarta at resolution~8 ($I = 0.030$, $p = 0.033$).

\medskip
This work is relevant to the FOSS4G community for four reasons. First, it demonstrates that PySAL, OSMnx, statsmodels, and H3---all open-source---can handle a dataset of 316 million observations end-to-end in approximately 15 minutes on consumer hardware (Mac Mini M2 Pro), establishing that FOSS4G tools are ready for operational, city-scale traffic analysis, not just prototyping. Second, the pipeline is released as a versioned, installable PyPI package with eleven CLI commands and comprehensive documentation, lowering the barrier for adoption by transportation agencies in resource-constrained settings where proprietary licenses are unaffordable. Third, the modular, command-per-analysis architecture provides a reusable template for applying FOSS4G tools to other urban analytics domains---air quality, land use change, public health surveillance---extending the impact beyond traffic. Fourth, the substantive finding---that congestion is primarily a demand synchronization problem rather than a spatial infrastructure constraint---carries direct policy implications for cities investing in road expansion versus demand management, and was only discoverable through the multi-method integration that Python's composable ecosystem uniquely enables.

The convergence of four independent spatial tests on null results for absolute speed, validated through MAUP-robust H3 aggregation, provides a level of methodological rigor that strengthens confidence in this conclusion. We believe this combination of open-source tooling, reproducible packaging, large-scale empirical validation, and policy-relevant findings makes a strong case for presentation at the FOSS4G Academic Track, and we look forward to engaging with the community on extending the pipeline to cities beyond Indonesia.

\medskip
\noindent\textbf{Keywords:} open-source, reproducibility, PySAL, OSMnx, traffic congestion,  FOSS4G

\end{document}